%=====================================================================
% FRONT MATTER
%=====================================================================
\pagenumbering{roman}
% Front-matter figures sit outside any chapter, so give them their own prefix
% rather than letting them inherit a misleading chapter number.
\structuralfigures{F.}
\thispagestyle{empty}

\begin{center}
\vspace*{1.0cm}

{\color{secondary}\rule{0.85\textwidth}{2pt}}\\[6pt]
{\large\textsc{\color{secondary}Bachelor of Science in Information Technology}}\\[18pt]

{\Huge\bfseries\color{primary}VIRTUAL SYSTEMS AND SERVICES}\\[10pt]
{\LARGE\bfseries\color{primary}From Hypervisors to the Cloud-Native Stack}\\[16pt]

{\color{secondary}\rule{0.85\textwidth}{2pt}}\\[16pt]

{\Large\textsc{Unified Lecture Handout}}\\[6pt]
{\large\color{gray!70!black}Complete Course in One Volume $\bullet$ Fall 2026}\\[24pt]

% --- the one claim the whole book rests on ---
\begin{tikzpicture}
  \node[rounded corners=4pt, draw=primary, line width=1.1pt, fill=primary!6,
        text=primary, align=center, font=\small\bfseries,
        minimum width=11.4cm, minimum height=1.5cm]
       {\faIcon{layer-group}\\[2pt]One idea, six times over: put a layer
        underneath, and lie to what sits above it};
\end{tikzpicture}\\[10pt]
{\small\itshape\color{gray!70!black}A CPU that is not a CPU, a disk that is not a
disk, a network that is not a network --- and the cost of each lie.}

\vspace{1.0cm}

\begin{tcolorbox}[enhanced, width=0.86\textwidth, colback=lightbg,
  colframe=primary, boxrule=0pt, leftrule=4pt, arc=2pt, top=6pt, bottom=6pt]
\small
\textbf{\color{primary}What this volume covers.} The HEC course \emph{Virtual
Systems and Services}, modernised. Parts~I and~II teach the 2017 outline close
to its published specification --- what a hypervisor is, why x86 resisted
virtualization for twenty years, and what Intel and AMD added to hardware to
end that. Parts~III to~VI carry the same ideas forward into containers,
Kubernetes, storage and network virtualization, cloud service models, live
migration, and confidential computing.

\smallskip
HEC published contents for this course once, in 2017. The 2025 revision drops
it and divides its ground between \emph{Cloud Computing} and \emph{Microservices
Architecture and Docker Containers}. Appendix~F states exactly where the line
falls, and why a course that goes one layer below both still earns its place.
\end{tcolorbox}

\vspace{0.9cm}
{\small\color{gray!60!black}Six Parts $\bullet$ Twenty-Three Chapters $\bullet$ Seven Appendices}
\vspace{1cm}
\end{center}

\newpage

%=====================================================================
\chapter*{How to Use This Handout}
\addcontentsline{toc}{chapter}{How to Use This Handout}
\markboth{How to Use This Handout}{}

This book has one subject, met six times. A \term{virtualization layer} sits
between something real and something that believes it has the real thing, and
keeps the belief intact. Chapters~3 to~7 do it to a processor. Chapters~8
to~11 do it to an operating system. Chapter~12 does it to a disk, Chapter~13
to a network, Chapter~15 to a desktop, Chapter~20 to a graphics card. The
machinery differs; the question never does --- \emph{what is the lie, who
maintains it, and what does maintaining it cost?}

Read it in order. Each layer is built from the one before.

\section*{The six parts}

\textbf{Part~I --- Foundations of Virtualization} (Chapters 1--3) fixes the
vocabulary and then proves the hard part. Chapter~3 is the theoretical spine of
the book: Popek and Goldberg's conditions, why x86 failed them for two decades,
and the two ways round that failure. Everything in Part~II is a hardware answer
to a problem stated in Chapter~3.

\textbf{Part~II --- Hypervisors and Compute Virtualization} (Chapters 4--7) is
the engine room. What Intel and AMD put in silicon; the four platforms you will
actually meet; how a CPU and a gigabyte of memory are divided among guests that
each think they own everything; and how a device is shared, emulated or handed
over outright.

\textbf{Part~III --- Containers and Lightweight Virtualization} (Chapters 8--11)
changes what is virtualized. Not the hardware, but the operating system's own
view of itself: namespaces, control groups and layered filesystems. Then Docker,
then Kubernetes, then the designs that put a hypervisor back underneath a
container because isolation mattered more than density.

\textbf{Part~IV --- Storage, Network and Cloud Virtualization} (Chapters 12--15)
applies the same idea to everything that is not a CPU. It ends at the cloud,
which is this book's six layers sold by the hour.

\textbf{Part~V --- Management, Performance and Security} (Chapters 16--19) is
the operations half: what you run to manage a fleet, how a running machine moves
between hosts without stopping, which numbers tell you a host is overloaded ---
and the attacks that exist only because the layer exists.

\textbf{Part~VI --- Emerging Hardware and Practice} (Chapters 20--23) closes
with the parts still moving: accelerator partitioning, encrypted memory and
attestation, four case studies, and your project.

\begin{conceptbox}
Why the theory chapter comes third, and not later. Chapter~3 is the most
demanding chapter in the book, and it is deliberately early. Without it,
hardware-assisted virtualization is a list of acronyms to memorise; with it,
every feature in Chapter~4 has a reason --- each one exists because
trap-and-emulate could not be made to work on x86. A reader who skips
Chapter~3 can still pass the course and will never understand why any of it
was built.
\end{conceptbox}

\section*{The furniture of every chapter}

Every chapter has the same parts, in the same order. Learn to read them once.

\rowcolors{2}{white}{lightbg}
\begin{longtable}{@{}L{4.6cm}L{10.6cm}@{}}
\toprule
\rowcolor{primary!15}
\textbf{Box} & \textbf{What it is for} \\
\midrule
\endhead
\faIcon{bullseye}~\textbf{Outcomes} & What you will be able to \emph{do}. Bloom
verbs, not topics. Use it as a checklist before the examination. \\
\faIcon{link}~\textbf{Before you start} & Which earlier chapters this one stands
on. If these are shaky, go back; nothing here is improved by pressing on. \\
\faIcon{tags}~\textbf{Key terms} & The vocabulary introduced. Every one appears
in the glossary of Appendix~D. \\
\faIcon{book}~\textbf{Definition} & The canonical statement. Learn these
verbatim; this field is full of terms used loosely in marketing and precisely in
documentation. \\
\faIcon{lightbulb}~\textbf{Concept} & The judgement call --- what an experienced
administrator would decide here, and why. \\
\faIcon{exclamation-triangle}~\textbf{Warning} & Something that will bite you,
usually in production and usually at scale. \\
\faIcon{calculator}~\textbf{Worked example} & Real arithmetic, every step shown.
Consolidation ratios, page-walk costs, migration times, availability budgets.
Do them with a pen. \\
\faIcon{exclamation-circle}~\textbf{Common pitfalls} & The mistakes that are
actually made, with the symptom each produces. \\
\faIcon{flask}~\textbf{Try it yourself} & The laboratory. Each is a runnable file
in \texttt{code/}, named for its chapter. \\
\faIcon{question-circle}~\textbf{Checkpoint} & Three questions. Attempt them
before reading the answers in Appendix~C. \\
\faIcon{clipboard-check}~\textbf{Chapter summary} & What to retain. \\
\faIcon{pen-fancy}~\textbf{Review questions} & Examination practice: multiple
choice, short answer, applied. \\
\bottomrule
\end{longtable}

\section*{How to study with it}

\begin{enumerate}[leftmargin=2.2em, itemsep=3pt, topsep=4pt]
  \item Read the outcomes first. They tell you what the chapter is \emph{for}.
  \item Read the prose once without stopping, then again with a pen on the
        worked examples. Every number in this book was computed before it was
        typed; you should be able to reproduce each one.
  \item Run the lab. This is not an optional enrichment: a hypervisor is a thing
        you develop intuition about by breaking it, and the labs are built to be
        broken. Overcommit memory until the guest swaps. Fill a thin pool.
        Migrate a machine while pinging it.
  \item Do the checkpoint from memory. Then read Appendix~C and argue with it.
  \item Attempt the applied review questions in pairs. They are written to be
        discussed, not looked up.
\end{enumerate}

\begin{alertbox}[On the laboratories, and what they need]
Every lab in this book runs on one student laptop with free software. Where the
subject is an enterprise platform --- vSphere, Hyper-V on a cluster, a Ceph
estate, an NVIDIA vGPU --- the chapter teaches the real thing and the lab uses a
free equivalent that behaves the same way, naming the substitution in its first
line. KVM with \texttt{virsh} stands in for vSphere; \texttt{kind} stands in for
a production Kubernetes cluster; LVM thin pools stand in for enterprise storage.

\smallskip
Three labs need hardware not every reader has: nested virtualization in
Chapter~4, GPU partitioning in Chapter~20, and an SEV- or TDX-capable processor
in Chapter~21. Each ships captured output in \texttt{code/} so the exercise can
be done by reading. No lab in this book is an instruction to buy something.

\smallskip
Appendix~A installs and verifies the whole stack in one sitting.
\end{alertbox}

\newpage

%=====================================================================
\chapter*{Course Specification}
\addcontentsline{toc}{chapter}{Course Specification}
\markboth{Course Specification}{}

This page is the course file. It states what HEC requires of a course called
\emph{Virtual Systems and Services}, and shows where in this handout each
requirement is met.

This course's position in the HEC curricula is unusual enough to state plainly
before the tables begin. HEC published contents for it \emph{once}, in the 2017
\emph{Curriculum of BS \& MS CS, SE \& IT} (pp.~71--130). The 2023 revision
keeps the title in the BSIT tables as a domain elective but publishes no
contents for the Information Technology cluster at all. The 2025 revision
removes the course, dividing its ground between \emph{Cloud Computing}, a new
major core, and \emph{Microservices Architecture and Docker Containers}.

The 2017 outline is therefore the only authoritative description that exists,
and this handout treats it as binding on Parts~I and~II and as a starting point
for the rest.

\section*{Course particulars}

\rowcolors{2}{white}{lightbg}
\begin{longtable}{@{}L{4.2cm}L{11.0cm}@{}}
\toprule
\rowcolor{primary!15}
\textbf{Item} & \textbf{Value} \\
\midrule
\endhead
Course title & Virtual Systems and Services \\
Programme & BS Information Technology. A domain elective under HEC 2017 and
HEC 2023; absent from HEC 2025 \\
Credit hours & 3 (2--1): two credit hours of lectures and one of laboratory \\
Prerequisite & Programming Fundamentals (HEC 2017). In practice this course also
assumes comfort at a Linux command line; the self-check two pages on says what
that means concretely \\
Duration & Sixteen teaching weeks, with a midterm in Week~8 and a project due in
Week~16 (Appendix~G) \\
Course introduction & The HEC 2017 aim, verbatim: ``This course will investigate
the current state of virtualization in computing systems. Virtualization at both
the hardware and software levels will be examined, with emphasis on the
hypervisor configurations of systems such as Vmware, Zen and Hyper-V. The
features and limitations of virtual environments will be considered, along with
several case studies used to demonstrate the configuration and management of
such systems. Para-virtualized software components will be analyzed and their
pros and cons discussed. Processor and peripheral support for virtualization
will also be examined, with a focus on emerging hardware features and the future
of virtualization.'' \\
Set texts & See the note below. The two titles HEC 2017 lists are not about this
subject; Appendix~E gives the references this handout actually uses \\
\bottomrule
\end{longtable}

\begin{alertbox}[The set texts in the published outline are the wrong books]
HEC 2017 lists two texts for this course: K.\ M.\ Stanney, \emph{Handbook of
Virtual Environments: Design, Implementation and Applications}, and G.\ Burdea,
\emph{Virtual Reality Technology}.

\smallskip
Both are virtual-\emph{reality} references --- head-mounted displays, haptics,
presence and simulator sickness. Neither mentions a hypervisor. The outline's own
description is unambiguously about system virtualization, so the reading list
appears to have been assembled from the course title rather than its contents.
This handout does not use them. Appendix~E lists the primary sources it does
use, beginning with Popek and Goldberg's 1974 paper, the Intel and AMD
architecture manuals, and the Xen, KVM and Firecracker papers.

\smallskip
The outline also spells Xen as ``Zen''. It is Xen, after the Xenoserver project
at Cambridge.
\end{alertbox}

\section*{Course learning outcomes}

HEC 2017 publishes an outline but no Bloom-levelled outcomes, and the later
editions publish neither. The six below are written against the 2017 outline and
traced to the chapters that deliver them.

\rowcolors{2}{white}{lightbg}
\begin{longtable}{@{}L{1.3cm}L{7.6cm}L{2.1cm}L{3.8cm}@{}}
\toprule
\rowcolor{primary!15}
\textbf{CLO} & \textbf{By the end of the course the student will be able to} &
\textbf{Bloom} & \textbf{Chapters} \\
\midrule
\endhead
CLO-1 & Explain what a virtualization layer is, and classify a given system by
the layer it virtualizes and the technique it uses & C2 Understand & 1, 2, 8 \\
CLO-2 & State the Popek--Goldberg conditions and explain why x86 failed them,
and how binary translation and hardware assistance each answer that failure &
C3 Apply & 3, 4 \\
CLO-3 & Configure and manage virtual machines and containers on at least three
platforms, including storage and networking & C3 Apply & 5--7, 9--13, 16 \\
CLO-4 & Analyze the performance of a virtualized workload, identify the
bottleneck from hypervisor metrics, and size a host for a given set of guests &
C4 Analyze & 6, 17, 18 \\
CLO-5 & Evaluate the security and the cost of a virtualization or cloud design,
and justify a deployment model against stated requirements & C5 Evaluate &
14, 19, 21, 22 \\
CLO-6 & Build, deploy, measure and document a working virtualized service &
C6 Create & 23, and every lab \\
\bottomrule
\end{longtable}

\section*{Coverage of the HEC 2017 course outline}

The published outline is a single descriptive paragraph rather than a topic
list. Each of its eleven clauses is traced below, in the order HEC writes them.
Nothing is omitted.

\rowcolors{2}{white}{lightbg}
\begin{longtable}{@{}L{6.4cm}L{2.0cm}L{6.6cm}@{}}
\toprule
\rowcolor{primary!15}
\textbf{HEC 2017 clause} & \textbf{Chapter} & \textbf{Where exactly} \\
\midrule
\endhead
``The current state of virtualization in computing systems'' & 1 & What the
layer is, the four motives for inserting one, and the line from CP-67 to the
cloud \\
``Virtualization at the hardware level'' & 4 & VT-x and AMD-V, the VMCS,
extended page tables, the IOMMU \\
``\ldots\ and software levels'' & 2, 3 & Full virtualization, paravirtualization
and binary translation; the virtual machine monitor itself \\
``Hypervisor configurations of systems such as VMware, Xen and Hyper-V'' & 5 &
All three, plus KVM, each as an architecture diagram and a comparison table \\
``The features \ldots\ of virtual environments'' & 6, 7, 17 & Overcommit,
ballooning, page sharing, NUMA; virtio and passthrough; snapshots and live
migration \\
``\ldots\ and limitations'' & 6, 18, 19 & What overcommit costs, how a
virtualized benchmark lies, and the attack surface the layer creates \\
``Several case studies \ldots\ configuration and management of such systems'' &
16, 22 & The management platforms in Chapter~16; four full case studies in
Chapter~22 \\
``Para-virtualized software components will be analyzed'' & 7 & The split-driver
model, virtqueues, and the PV driver sets of each platform \\
``\ldots\ and their pros and cons discussed'' & 7 & Emulated against
paravirtualized against passthrough, measured, with the cost of each \\
``Processor \ldots\ support for virtualization'' & 4 & Chapter~4 entire \\
``\ldots\ and peripheral support'' & 7, 20 & SR-IOV and VFIO in Chapter~7; GPU
and accelerator partitioning in Chapter~20 \\
``Emerging hardware features and the future of virtualization'' & 20, 21 &
Accelerator virtualization; confidential computing, attestation, CXL and
composable infrastructure \\
\bottomrule
\end{longtable}

\section*{Topics taught here that the 2017 outline does not name}

The 2017 outline predates Docker's general adoption, Kubernetes entirely, and
every confidential-computing platform in Chapter~21. Eight chapters go beyond
its letter. They are listed here so that the departure is explicit.

\rowcolors{2}{white}{lightbg}
\begin{longtable}{@{}L{3.1cm}L{8.2cm}L{3.7cm}@{}}
\toprule
\rowcolor{primary!15}
\textbf{Chapters} & \textbf{Why they are here} & \textbf{What they rest on} \\
\midrule
\endhead
8--11 Containers & The dominant form of virtualization in industry since 2015,
and the direct subject of the course HEC 2025 created to replace this one.
Taught as what they are --- operating-system-level virtualization, one layer
above the hypervisor & HEC 2025, \emph{Microservices Architecture and Docker
Containers} \\
12--13 Storage and network & The 2017 clause ``virtualization at both the
hardware and software levels'' is not restricted to the processor, and no
hypervisor is useful without virtual disks and virtual switches & HEC 2017,
``hardware and software levels'' \\
14 Cloud models & Where virtualization is actually consumed. Without it the
course ends at the machine room door & HEC 2025, \emph{Cloud Computing} \\
21 Confidential computing & The direct continuation of ``emerging hardware
features and the future of virtualization'', written sixteen years later &
HEC 2017, ``emerging hardware features'' \\
\bottomrule
\end{longtable}

\begin{conceptbox}
The overlap with \emph{Cloud Computing} is real, intended, and bounded. The
curriculum map that traces these editions reaches a blunt conclusion ---
``virtualisation is one chapter of cloud computing, not a peer of it'' --- and
where a department offers both courses, that is the right call.

\smallskip
Where this one is offered, it earns its three credit hours by going a layer
\emph{below} a cloud course rather than alongside it. A cloud course explains
what an instance is and how to rent one. This course explains what happens on
the host when you do: which instruction trapped, which page table was walked
twice, which driver was split in half, and why the instance beside yours can
make yours slow. Appendix~F draws the line topic by topic, and names what this
course deliberately does not teach.
\end{conceptbox}

\newpage

%=====================================================================
\chapter*{The Course at a Glance}
\addcontentsline{toc}{chapter}{The Course at a Glance}
\markboth{The Course at a Glance}{}

\dsfigh{%
\begin{tikzpicture}[
  part/.style={rounded corners=3pt, fill=#1, text=white, font=\bfseries\footnotesize,
               minimum width=2.5cm, text width=2.2cm, minimum height=0.95cm,
               align=center},
  ch/.style={rounded corners=2pt, draw=#1, line width=0.7pt, fill=#1!7, text=#1!60!black,
             font=\scriptsize, minimum width=2.5cm, minimum height=0.52cm, align=center},
  arr/.style={-{Stealth[length=2mm]}, primary!45, line width=0.8pt}
]
\node[part=primary]   (p1) at (0,0)    {Part I\\Foundations};
\node[part=secondary] (p2) at (3.0,0)  {Part II\\Hypervisors};
\node[part=greenhl]   (p3) at (6.0,0)  {Part III\\Containers};
\node[part=purplehl]  (p4) at (9.0,0)  {Part IV\\Storage, Network, Cloud};
\node[part=tealhl]    (p5) at (12.0,0) {Part V\\Operations};
\node[part=goldhl]    (p6) at (15.0,0) {Part VI\\Emerging and Practice};

\foreach \a/\b in {p1/p2,p2/p3,p3/p4,p4/p5,p5/p6}{\draw[arr] (\a) -- (\b);}

\foreach \i/\t in {1/{1. What Virtualization Is}, 2/{2. Architectures},
                   3/{3. The VMM}}{
  \pgfmathsetmacro{\yy}{-0.90-(\i-1)*0.62}
  \node[ch=primary] at (0,\yy) {\t};}

\foreach \i/\t in {1/{4. Hardware Support}, 2/{5. Platforms},
                   3/{6. CPU and Memory}, 4/{7. Paravirtualization}}{
  \pgfmathsetmacro{\yy}{-0.90-(\i-1)*0.62}
  \node[ch=secondary] at (3.0,\yy) {\t};}

\foreach \i/\t in {1/{8. VMs to Containers}, 2/{9. Docker},
                   3/{10. Kubernetes}, 4/{11. Lightweight}}{
  \pgfmathsetmacro{\yy}{-0.90-(\i-1)*0.62}
  \node[ch=greenhl] at (6.0,\yy) {\t};}

\foreach \i/\t in {1/{12. Storage}, 2/{13. Network and SDN},
                   3/{14. Cloud Models}, 4/{15. Desktop and App}}{
  \pgfmathsetmacro{\yy}{-0.90-(\i-1)*0.62}
  \node[ch=purplehl] at (9.0,\yy) {\t};}

\foreach \i/\t in {1/{16. Management}, 2/{17. HA and Migration},
                   3/{18. Performance}, 4/{19. Security}}{
  \pgfmathsetmacro{\yy}{-0.90-(\i-1)*0.62}
  \node[ch=tealhl] at (12.0,\yy) {\t};}

\foreach \i/\t in {1/{20. GPU and Accelerators}, 2/{21. Confidential Computing},
                   3/{22. Case Studies}, 4/{23. Project}}{
  \pgfmathsetmacro{\yy}{-0.90-(\i-1)*0.62}
  \node[ch=goldhl] at (15.0,\yy) {\t};}

\node[rounded corners=3pt, fill=primary!7, draw=primary!30,
      minimum width=17.3cm, minimum height=0.66cm,
      font=\scriptsize\bfseries, text=primary] at (7.5,-3.95)
      {\faIcon{layer-group}~~The same question in every part: what is the lie,
       who maintains it, and what does maintaining it cost?};
\end{tikzpicture}}%
{The six parts and the chapters in each. Parts I and II teach the HEC 2017
outline close to specification; Parts III to VI carry it forward to the stack
in production use today.}{coursemap}

\vspace{4pt}
{\small\itshape\color{gray!70!black}
This course sits beside \emph{Cloud Computing}, \emph{System and Network
Administration} and the CCNA track, and one layer below all three.
Appendix~F states what belongs to each.}

\newpage

%=====================================================================
\chapter*{The Laboratory Stack}
\addcontentsline{toc}{chapter}{The Laboratory Stack}
\markboth{The Laboratory Stack}{}

This course is unusual among the BSIT electives in that the subject is
infrastructure, and infrastructure is expensive. A faithful laboratory for the
2017 outline would need a vSphere licence, a Hyper-V cluster, a SAN and a
GPU with vGPU entitlement --- perhaps forty thousand dollars of software before
any hardware.

That is not what this book asks for. Every lab runs on one laptop with free
software, and the substitution is named each time it is made.

\rowcolors{2}{white}{lightbg}
\begin{longtable}{@{}L{4.2cm}L{5.2cm}L{5.6cm}@{}}
\toprule
\rowcolor{primary!15}
\textbf{What the chapter teaches} & \textbf{What the lab uses} &
\textbf{Why the substitution is honest} \\
\midrule
\endhead
VMware ESXi and vSphere & KVM with \texttt{virsh} and \texttt{virt-manager} &
Both are Type-1 hypervisors with the same primitives. The vocabulary differs;
the mechanism does not \\
Microsoft Hyper-V & Hyper-V where Windows Pro is available, KVM otherwise &
Hyper-V is free with the operating system most students already run \\
vMotion / Live Migration & \texttt{virsh migrate} between two nested hosts &
Pre-copy migration with dirty-page tracking, which is the mechanism
Chapter~17 teaches \\
Enterprise SAN, Ceph, vSAN & LVM thin pools; a single-node \texttt{cephadm}
cluster & Thin provisioning, copy-on-write and the over-subscription failure are
all reproducible on one disk \\
Production Kubernetes & \texttt{kind} (Kubernetes in Docker) & A real API
server, scheduler and kubelet. Only the node count is a pretence \\
NVIDIA vGPU and MIG & \texttt{nvidia-smi} partitioning where a GPU exists;
captured output otherwise & Partitioning policy is the lesson; the silicon is
not available to rent by the hour \\
AWS, Azure and GCP & The three free tiers, with a teardown step in every lab &
The instance you launch is the instance the chapter describes \\
Intel TDX, AMD SEV & Captured attestation output in \texttt{code/} & Confidential
VMs need a specific processor generation and a cooperating host \\
\bottomrule
\end{longtable}

\begin{alertbox}[One thing the free stack cannot teach you]
Scale. Every pathology in Chapter~18 --- the noisy neighbour, the oversubscribed
cluster, the storage array that saturates at three in the afternoon --- is a
property of many machines sharing one resource, and a laptop running four guests
will not show it to you honestly.

\smallskip
Where that is the case the chapter says so and gives measured numbers from real
estates instead of asking you to produce them. Read those numbers as data, and
treat the lab as the mechanism rather than the magnitude.
\end{alertbox}

Appendix~A installs the whole stack and verifies it. Appendix~B is the command
reference you will keep open beside the labs.

\newpage

%=====================================================================
\chapter*{Notation and Conventions}
\addcontentsline{toc}{chapter}{Notation and Conventions}
\markboth{Notation and Conventions}{}

\rowcolors{2}{white}{lightbg}
\begin{longtable}{@{}L{3.4cm}L{11.8cm}@{}}
\toprule
\rowcolor{primary!15}
\textbf{Term or symbol} & \textbf{Meaning in this book} \\
\midrule
\endhead
\term{host} & The physical machine, and the software running directly on it \\
\term{guest} & An operating system running inside a virtual machine \\
\term{hypervisor}, \term{VMM} & Used interchangeably, as the literature does.
Where the distinction matters --- KVM is a VMM inside a host kernel, not a
free-standing hypervisor --- the text says which it means \\
Ring 0, 1, 3 & x86 privilege levels: kernel, the level Xen deprivileged guests
into, and user \\
\texttt{root}, \texttt{non-root} & VMX operation modes, Chapter~4. Nothing to do
with the \texttt{root} user \\
GVA, GPA, HPA & Guest virtual, guest physical and host physical address --- the
three address spaces a nested page walk crosses \\
vCPU & A virtual processor presented to a guest; scheduled onto a physical core
by the hypervisor \\
pCPU & A physical core or hardware thread \\
$R$ & Consolidation ratio: guests per host \\
$\sigma$ & Overcommit ratio: allocated resource divided by physical resource \\
\%RDY & Fraction of time a vCPU was runnable but not scheduled. The single most
useful number in Chapter~18 \\
RPO, RTO & Recovery point objective (how much data you may lose) and recovery
time objective (how long you may be down) \\
p95, p99 & The 95th and 99th percentile of a latency distribution \\
1 GiB, 1 GB & Binary and decimal gigabytes, $2^{30}$ and $10^{9}$ bytes. Storage
vendors mean the second; operating systems usually report the first. This book
is explicit every time it matters \\
\bottomrule
\end{longtable}

\begin{conceptbox}
Two conventions worth stating because the industry is inconsistent about them.

\smallskip
\textbf{Type-1 and Type-2} are used in this book as Goldberg defined them in
1973 --- by whether the VMM runs on bare metal or on a host operating system ---
and Chapter~2 shows why the distinction is blurrier in practice than the
textbook diagram suggests. KVM is the awkward case, and it is discussed rather
than assigned.

\smallskip
\textbf{Virtualization and emulation} are not synonyms. Virtualization runs
guest instructions on the real processor wherever it can; emulation interprets
every one of them. QEMU does both, in different modes, which is why it appears
on both sides of the comparison in Chapter~2.
\end{conceptbox}

\newpage

%=====================================================================
\chapter*{Prerequisite Self-Check}
\addcontentsline{toc}{chapter}{Prerequisite Self-Check}
\markboth{Prerequisite Self-Check}{}

The published prerequisite is \emph{Programming Fundamentals}, which is not
really what this course needs. What it needs is the twelve things below. If you
can do ten of them, begin at Chapter~1. If you cannot, Appendix~A and the first
two sections of Appendix~B will close most of the gap in an afternoon.

\begin{enumerate}[leftmargin=2.4em, itemsep=4pt, topsep=4pt]
  \item Explain what an operating system kernel does that a user program cannot.
        \hfill{\scriptsize\color{neutral}(Operating Systems)}
  \item Describe what happens on a system call, from the user instruction to the
        return. \hfill{\scriptsize\color{neutral}(\chref{vmm})}
  \item Explain virtual memory: what a page table maps, and what the MMU does
        with it. \hfill{\scriptsize\color{neutral}(\chref{hardware})}
  \item State what a process is, and what the kernel keeps for each one.
        \hfill{\scriptsize\color{neutral}(\chref{containers})}
  \item Navigate a Linux filesystem and read a file from the command line
        without a graphical tool. \hfill{\scriptsize\color{neutral}(Appendix~B)}
  \item Edit a text file over SSH and explain what SSH is doing.
        \hfill{\scriptsize\color{neutral}(Appendix~A)}
  \item Read a shell script containing a loop, a conditional and a pipeline.
        \hfill{\scriptsize\color{neutral}(Appendix~B)}
  \item Explain the difference between an IP address, a subnet mask and a
        default gateway. \hfill{\scriptsize\color{neutral}(\chref{network})}
  \item Say what a TCP port is and how two services share one host.
        \hfill{\scriptsize\color{neutral}(\chref{network})}
  \item Convert between binary, decimal and hexadecimal, and read a size in
        bytes written either way. \hfill{\scriptsize\color{neutral}(Notation,
        above)}
  \item Explain what a file format is, given that a virtual disk is a file that
        contains a filesystem. \hfill{\scriptsize\color{neutral}(\chref{storage})}
  \item Describe what a compiler produces and in what sense a CPU
        ``runs'' it. \hfill{\scriptsize\color{neutral}(Programming
        Fundamentals)}
\end{enumerate}

\begin{conceptbox}
Notice what is \emph{not} on that list: no calculus, no statistics, no
networking certification, and no prior exposure to any cloud platform. The
arithmetic in this book is ratios, percentages and powers of two. What it
genuinely demands instead is comfort with the idea that software can be
\emph{underneath} other software --- which is uncomfortable the first time, and
then never again.
\end{conceptbox}

\newpage
\tableofcontents
\newpage
\listoffigures
\newpage

\pagenumbering{arabic}
